% Read STYLE.md before editing prose in this file.
% Follow the Content Creation Workflow in WORKFLOW.md §Content Creation Workflow.

\section{Background}
\label{sec:background}

\subsection{Cetacean Acoustic Communication}
\label{subsec:cetacean-acoustics}

Bottlenose dolphins (\spname{Tursiops truncatus}) produce three broad
classes of sound: (1)~frequency-modulated whistles, which serve as
individual identification signals or ``signature whistles'';
(2)~burst-pulse sounds, a rapid series of clicks used in social
contexts; and (3)~click trains for echolocation~\citep{Wonko2021communication}.
Research into the semantic content of these vocalizations has accelerated
over the past two decades, yet leave-taking behavior has been studied
primarily from an ethological rather than an acoustic perspective.

\citet{Slartibartfast2023topology} note, almost in passing, that coastal
populations near fjord environments exhibit unusually structured departure
calls, attributing this to the particular resonant properties of such
coastlines.
This observation motivated our own data collection strategy, described
in \cref{sec:methodology}.

\subsection{Computational Bioacoustics}
\label{subsec:comp-bioacoustics}

Sequence-level analysis of animal vocalizations has benefited greatly
from methods developed for human speech processing.
Hidden Markov models~(HMMs) were first applied to bird song classification
by \citet{Deep2022number}, who demonstrated that temporal dependencies
within a vocalization bout could be captured by the Markov transition
matrix, achieving state-of-the-art accuracy on a 12-species corpus.

More recently, \citet{Trillian2023signal} applied Fourier-based spectral
envelope analysis to burst-pulse sequences, showing that departure-adjacent
bouts in captive dolphins exhibited measurably broader spectral envelopes
than those recorded during feeding events.
Our work extends this result to wild populations and tests whether the
pattern generalizes across geographically distinct groups.

\begin{figure}[htbp]
	\centering
	% TODO: Replace the placeholder below with the actual spectrogram image:
	%   figures/ch01/spectrogram_example.pdf
	\fbox{\rule{0pt}{4cm}\hspace{8cm}}%
	\caption{Representative spectrogram of an \aseq{extended-farewell}
		sequence recorded at Site~A (Atlantic, 2023). The vertical axis
		shows frequency in \kHz{}; the horizontal axis shows time in
		seconds. Color indicates power spectral density.
		(Placeholder: replace with \texttt{figures/ch01/spectrogram\_example.pdf}.)
	}
	\label{fig:ch01-spectrogram-example}
\end{figure}

\Cref{fig:ch01-spectrogram-example} illustrates a representative
\aseq{extended-farewell} sequence; note the characteristic frequency
sweep from approximately 8\kHz{} to 18\kHz{} occurring in the
final 1.2~seconds before departure.

\subsection{Leave-Taking in Animal Communication}
\label{subsec:leave-taking}

The study of leave-taking rituals in animal communication is a nascent
field~\citep{Marvin2022depression}.
Across species ranging from great apes to corvids, researchers have
documented the existence of structured terminal vocalizations that differ
acoustically from vocalizations produced in non-departure contexts.
\citet{Prefect2025syntax} provides the most direct antecedent to our
work: a minimalist syntactic analysis of 17~\spname{Tursiops} farewell
calls recorded near the Azores, arguing for a recursive hierarchical
structure in the calls' temporal organization.
Our study corroborates Prefect's syntactic findings while providing the
first corpus-scale statistical evidence for their generality.
